\section{Exogenous Claims in GFM}
\label{sec:gfm_claims}

% This section instantiates the general commitment protocol (§5) for
% the three kinds of claims that GFM agents make. Each claim type
% maps onto the protocol's (Commit, Verify, Reveal) phases with a
% claim-specific exercise protocol.

% §7.1 Capability claims
%   Claim: "I have capability c_k in the poset."
%   Commitment: Com(c_k, r) where c_k is a node in the capability
%     poset of Paper 2.
%   Exercise protocol: the substrate-exclusive witness tests whether
%     the agent can exercise c_k. "Exercise" is defined through
%     Paper 3's exercise-pathway formalism: the witness constructs a
%     scenario that requires c_k and observes whether the agent
%     succeeds.
%   Verification result: pass (capability confirmed) or fail
%     (commitment falsified). Falsified commitments produce a large
%     negative trust update (§8).
%
%   Key issue: what counts as a valid exercise protocol for
%   capability c_k? This is domain-specific. For concrete capabilities
%   (e.g., "can solve differential equations") the exercise protocol
%   is a test. For abstract capabilities (e.g., "strategic planning")
%   the exercise protocol may require behavioral observation over time.
%   The paper should define the interface (what an exercise protocol
%   must satisfy) without prescribing the implementation.

% §7.2 Risk claims
%   Claim: "Exercise pathway P has probability p and contraction
%     magnitude Δ."
%   Commitment: Com((P, p, Δ), r) using the pathway-conditional
%     risk formalism of Paper 4 Definition 3.
%   Exercise protocol: the witness evaluates the pathway using its
%     own SCM (Paper 4). The structural verification residual
%     r^risk_j = (p_claimed - p_witness, Δ_claimed - Δ_witness)
%     is computed by the witness, not the claimant.
%   Verification result: the residual norm. If ||r^risk||² exceeds
%     a threshold, the claim is flagged as unreliable.
%
%   Key issue: the witness's SCM may differ from the claimant's.
%   The commitment protocol does not resolve SCM disagreements; it
%   ensures that the claimant cannot modify the witness's evaluation.
%   SCM disagreements are handled through the multi-channel
%   attribution framework of Paper 4 §2.4.

% §7.3 Benchmark validations
%   Claim: "Agent a_j passes benchmark B with score s."
%   Commitment: Com((B, s), r).
%   Exercise protocol: the witness runs benchmark B against agent a_j
%     independently. The witness must have access to the benchmark
%     specification but not to the claimant's self-evaluation.
%   Verification result: the witness's independently measured score
%     s_witness. If |s - s_witness| exceeds a threshold, the
%     benchmark claim is flagged.
%
%   Key issue: coalition capture. A coalition can design benchmark B
%   such that only coalition members pass. The commitment protocol
%   resists this because the witness is substrate-exclusive: the
%   coalition cannot control the witness's execution of B.
%   But the coalition CAN design B itself to be biased. The defense
%   is that benchmark specifications are public (committed to the
%   ledger) and auditable by any agent. A biased benchmark is
%   detectable by third parties who can evaluate B's coverage.

% TODO: formal definitions mapping each claim type onto the
% general protocol, with claim-specific security properties.
